\section{Niveau adéquat de maturité (G2)}

La méthode note douze questions de 0 à 3, garde le maximum de chaque famille de trois et convertit la somme des quatre maxima en niveau. La figure~\ref{fig:notes} présente les notes et le calcul. Les justifications complètes, pièce par pièce, figurent en annexe~\ref{ann:autodiag}.

\begin{figure}[H]
\centering
\resizebox{\linewidth}{!}{% Notes Q01 à Q12, maxima par famille et calcul du niveau adéquat (onglet 01)
\begin{tikzpicture}[x=0.82cm, y=0.95cm, font=\scriptsize]
  % familles
  \foreach \fam/\kcol/\deb in {Conséquences/accent/0, Sensibilité/vert/3, Menaces/ambre/6, Vulnérabilités/violet/9}{
    \fill[\kcol!7] (\deb+0.1,-0.95) rectangle (\deb+2.9,3.35);
    \node[font=\scriptsize\bfseries, text=\kcol] at (\deb+1.5,-0.7) {\fam};}
  % grille
  \foreach \y in {0,1,2,3}{\draw[filet] (0,\y) -- (12,\y); \node[anchor=east, text=gris] at (-0.05,\y) {\y};}
  % barres
  \foreach \q/\n [count=\i from 0] in \NotesQ{
    \pgfmathtruncatemacro\fam{div(\i,3)}
    \ifcase\fam \def\kcol{accent}\or\def\kcol{vert}\or\def\kcol{ambre}\or\def\kcol{violet}\fi
    \fill[\kcol] (\i+0.25,0) rectangle (\i+0.75,\n);
    \node[above, font=\scriptsize\bfseries] at (\i+0.5,\n) {\n};
    \node[text=gris] at (\i+0.5,-0.25) {\q};}
  % maxima
  \foreach \m/\deb [count=\j] in {\MaxConseq/0, \MaxSensib/3, \MaxMenaces/6, \MaxVuln/9}{
    \draw[encre, thick, dashed] (\deb+0.15,\m) -- (\deb+2.85,\m);
    \node[font=\scriptsize, text=encre, fill=white, inner sep=1pt] at (\deb+1.5,3.15) {max $=\m$};}
  % calcul
  \node[draw=accent, fill=accentclair, rounded corners=2pt, align=center, font=\small, inner sep=6pt, anchor=west] at (12.5,1.4)
    {$\MaxConseq + \MaxSensib + \MaxMenaces + \MaxVuln = \textbf{\SommeMaxima}$\\[3pt]
     \scriptsize barème 2007 : 9 à 10 donne 4\\[3pt]
     \textbf{niveau adéquat \NiveauCalcule}\\[1pt]\scriptsize ajustement \Ajustement, cible \Cible};
\end{tikzpicture}
}
\caption{Notes Q01 à Q12, maxima par famille et niveau adéquat (onglet 01).}
\label{fig:notes}
\end{figure}

\begin{table}[H]
\centering\small
\caption{Autodiagnostic DCSSI, synthèse}
\begin{tabularx}{\linewidth}{@{}C{1cm} Y C{1cm} L{3.2cm} L{2.2cm}@{}}
\toprule
\entete \hd{ID} & \hd{Question} & \hd{Note} & \hd{Pièces} & \hd{Confiance} \\
\midrule
\lignes{tab_questions_court}
\bottomrule
\end{tabularx}
\end{table}

Le résultat est robuste. La cible resterait à 4 si la confidentialité était notée 2, la somme valant encore 9. Elle ne passerait à 5 que si l'ouverture du SI était jugée hors de contrôle à cause de l'accès public au portail. Nous ne retenons pas cette lecture, car la publication est unique, filtrée et protégée par un reverse proxy \piece{(T01, T02)}.

La méthode autorise un ajustement d'un niveau. Une baisse à 3 pourrait s'appuyer sur la taille de l'entreprise et sur un RSSI à mi-temps \piece{(P02)}. Nous ne la proposons pas, car les données touchent la santé de patients. La contrainte de moyens est mieux traitée par une progression par étapes. La décision revient à la DG.
